\documentclass[11pt,a4paper]{article}
\usepackage[margin=27mm]{geometry}
\usepackage[T1]{fontenc}
\usepackage{lmodern,amsmath,amssymb,amsthm,mathtools,braket,dsfont}
\usepackage[colorlinks=true,linkcolor=blue,urlcolor=blue,citecolor=blue]{hyperref}
\usepackage{microtype}
\DeclareMathOperator{\Tr}{Tr}
\DeclareMathOperator{\Sym}{Sym}
\newcommand{\A}{\mathcal A}
\newcommand{\id}{\mathds{1}}
\newcommand{\E}{\mathbb E}
\newtheorem{theorem}{Theorem}
\newtheorem{lemma}{Lemma}
\newtheorem{remark}{Remark}
\title{Fidelity total anticoherence is nondecreasing\\under covariant noise at order two}
\author{}
\date{13 November 2026}
\begin{document}
\maketitle
\begin{abstract}
We prove the fidelity part of Conjecture~1 of \emph{Total, quantum, and
classical measures of anticoherence for mixed spin states}, at $t=2$, for
every $N\ge3$. The notation and the squared Uhlmann--Jozsa fidelity
normalization are those of \texttt{resub\_v0.tex}. The proof is based on
\texttt{su2\_t2\_full\_direct\_proof.pdf}. It treats arbitrary covariant
channels directly, derives their two relevant multipliers without $6j$
symbols, and ends with an elementary three-eigenvalue argument.
\end{abstract}

\section{Statement and conventions}
As in the manuscript, $\mathcal H_j$ is the spin-$j$ Hilbert space,
$N=2j$, and $\rho_S$ is the state corresponding to
$\rho\in\mathcal D(\mathcal H_j)$ under the identification
$\mathcal H_j\simeq\Sym^N(\mathbb C^2)$. Put
\[
 \rho_t=\Tr_{N-t}(\rho_S),\qquad
 \rho_0^{(t)}=\frac{\id_{t+1}}{t+1},\qquad
 F(\rho_t,\rho_0^{(t)})=\frac{(\Tr\sqrt{\rho_t})^2}{t+1}.
\]
Thus the fidelity total candidate of the manuscript is
\begin{equation}\label{eq:candidate}
 \A_t^{T,F}(\rho)=\frac{(\Tr\sqrt{\rho_t})^2-1}{t}.
\end{equation}
All determinants and traces of two-qubit reductions below are taken on
the three-dimensional symmetric subspace, including any zero eigenvalues.

\begin{theorem}[Conjecture 1, fidelity candidate, $t=2$]\label{thm:main}
Let $N\ge3$, $j=N/2$, and let $\Phi$ be a completely positive,
trace-preserving map on $\mathcal B(\mathcal H_j)$ satisfying
\[
 \Phi(U\rho U^\dagger)=U\Phi(\rho)U^\dagger
 \quad\text{for every spin-$j$ rotation $U$ and every state $\rho$.}
\]
Then, for every $\rho\in\mathcal D(\mathcal H_j)$,
\begin{equation}\label{eq:main}
 \boxed{\A_2^{T,F}\bigl(\Phi(\rho)\bigr)\ge\A_2^{T,F}(\rho).}
\end{equation}
\end{theorem}
The condition $N\ge3$ is exactly the proper-order condition $2\le t<N$
at $t=2$. No restriction to pure inputs or to random rotations is made.
The theorem does not address the trace-distance candidate.

\section{The two channel multipliers}
Use the Hilbert--Schmidt orthonormal irreducible tensor operators
$T_{LM}$ of the manuscript. The conjugation representation on
$\mathcal B(\mathcal H_j)$ contains one copy of each rank $L=0,\ldots,N$.
Schur's lemma and Hermiticity preservation therefore give real numbers
$f_L$ such that
\[
 \Phi(T_{LM})=f_LT_{LM}.
\]
Covariance implies that $\Phi(\id)$ is a scalar multiple of $\id$;
trace preservation fixes that scalar to one. In particular $f_0=1$.
Write
\[
 f=f_1,\qquad g=f_2,\qquad S=N(N+2),\qquad D=S(S-3)>0.
\]
The symbol $S$ here is a scalar, distinct from the subscript in $\rho_S$.

\begin{lemma}[Multiplier bounds]\label{lem:multipliers}
Every channel in Theorem~\ref{thm:main} satisfies
\begin{align}
 -\frac35\le g\le1,\qquad 3|f|\le2+g,\label{eq:region}\\
 3Sf^2-6f-S\le2(S-3)g.\label{eq:curved}
\end{align}
Consequently $|f|\le1$ and $|g|\le1$; if $g<-1/2$, then $|f|<1/3$.
\end{lemma}

\begin{proof}
We first derive formulas that include all covariant channels. Set
$d=N+1$ and $c=j(j+1)=S/4$. Let $C_\Phi$ be the unnormalized Choi
matrix, with convention
\begin{equation}\label{eq:choi-trace}
 \Tr[C_\Phi(B\otimes C^{\mathsf T})]=\Tr[B\Phi(C)].
\end{equation}
The unitary $V\ket{j,m}=(-1)^{j-m}\ket{j,-m}$ obeys
$VJ_a^{\mathsf T}V^\dagger=-J_a$. Consequently
$\widetilde C_\Phi=(\id\otimes V)C_\Phi(\id\otimes V^\dagger)$
commutes with $U\otimes U$. In $\mathcal H_j\otimes\mathcal H_j$,
weight $m\in\{-N,\ldots,N\}$ has multiplicity $N+1-|m|$.
Subtracting consecutive nonnegative-weight multiplicities gives one
copy of each spin $K=0,\ldots,N$. Thus the decomposition
$\mathcal H_j\otimes\mathcal H_j=\bigoplus_{K=0}^N\mathcal H_K$
is multiplicity free. Positivity of the Choi matrix and
$\Tr C_\Phi=d$ imply
\begin{equation}\label{eq:choi}
 \widetilde C_\Phi=d\sum_{K=0}^N p_K\frac{P_K}{2K+1},\qquad
 p_K\ge0,\qquad\sum_{K=0}^N p_K=1,
\end{equation}
where $P_K$ projects onto total spin $K$. We use $\E$ for this finite
weighted average and put $\chi_K=K(K+1)$.

For completeness the required Choi moments can be computed using only
angular-momentum commutators. Define
\[
 X=\sum_{a=1}^3J_a\otimes J_a,\qquad
 Q_{ab}=\frac{J_aJ_b+J_bJ_a}{2}-\frac c3\delta_{ab}\id.
\]
The total-spin Casimir gives $X P_K=(\chi_K/2-c)P_K$. The $J_a$
span rank one, whereas the symmetric traceless $Q_{ab}$ span rank two.
Moreover
\begin{equation}\label{eq:quadrupole}
 \sum_{a,b}Q_{ab}\otimes Q_{ab}
 =X^2+\frac X2-\frac{c^2}3\id,
 \qquad
 \sum_{a,b}Q_{ab}^2=\frac{c(4c-3)}6\id.
\end{equation}
Here the first identity follows by expanding the symmetrized products
and using $[J_a,J_b]=i\sum_k\epsilon_{abk}J_k$. For the second,
expansion of $\sum_b[J_b,[J_b,J_a]]=2J_a$ gives
$\sum_bJ_bJ_aJ_b=(c-1)J_a$. Thus, writing
$B_{ab}=(J_aJ_b+J_bJ_a)/2$,
\[
 \sum_{a,b}B_{ab}^2
 =\frac12\sum_{a,b}(J_aJ_bJ_aJ_b+J_aJ_bJ_bJ_a)
 =\left(c^2-\frac c2\right)\id,
\]
which proves the second identity in \eqref{eq:quadrupole}.

Time reversal changes the sign of $J_a^{\mathsf T}$ and fixes
$Q_{ab}^{\mathsf T}$. Applying \eqref{eq:choi-trace} therefore yields
\[
 -dcf=\Tr(\widetilde C_\Phi X)=d\E(\chi/2-c),
\]
and
\[
 \frac{dc(4c-3)}6g
 =d\E\left[(\chi/2-c)^2+\frac{\chi/2-c}{2}-\frac{c^2}3\right].
\]
Equivalently,
\begin{equation}\label{eq:fg}
 f=\E f_K,\quad g=\E g_K,\qquad
 f_K=1-\frac{2\chi_K}{S},\quad
 g_K=1+\frac{6\chi_K(\chi_K-S+1)}D.
\end{equation}
Since $0\le\chi_K\le N(N+1)<S-1$, we have $g_K\le1$.
For $N=3$, the four values of $g_K$ are $1,1/5,-3/5,1/5$.
For $N\ge4$, $S\ge24$, completing the square gives
\[
 g_K\ge1-\frac{3(S-1)^2}{2S(S-3)}\ge-\frac35,
\]
where the last inequality follows from
\[
 16S(S-3)-15(S-1)^2=(S-24)(S+6)+129>0.
\]
The two linear bounds follow on averaging the identities
\[
 2+g_K-3f_K=\frac{6\chi_K(\chi_K-2)}D\ge0,
 \qquad
 2+g_K+3f_K=\frac{6[(\chi_K-S+2)^2+S-4]}D>0.
\]
The first inequality uses the discrete fact that $\chi_K=0$ or
$\chi_K\ge2$. Finally, direct substitution into \eqref{eq:fg} gives
\begin{equation}\label{eq:variance}
 2(S-3)g-(3Sf^2-6f-S)
 =\frac{12}{S}\E[(\chi-\E\chi)^2]\ge0.
\end{equation}
If $g<-1/2$, \eqref{eq:curved} implies $Sf^2-2f-1<0$.
Since $S\ge15$, this forces $-1/3<f<1/3$: otherwise
$Sf^2-2f-1\ge15f^2-2f-1=(3f-1)(5f+1)\ge0$.
\end{proof}

\section{Physical two-qubit reductions}
Write the symmetric two-qubit reduction in its Pauli form
\begin{equation}\label{eq:pauli}
 \rho_2=\frac14\left[\id\otimes\id+
 \sum_{a=1}^3r_a(\sigma_a\otimes\id+\id\otimes\sigma_a)
 +\sum_{a,b=1}^3T_{ab}\sigma_a\otimes\sigma_b\right].
\end{equation}
Here $\mathbf r\in\mathbb R^3$, $T=T^{\mathsf T}\in M_3(\mathbb R)$,
and $\Tr T=1$, since $\sum_a\sigma_a\otimes\sigma_a=2\mathsf F-\id$
and the swap $\mathsf F$ is the identity on the symmetric subspace.
In the orthonormal Cartesian spin-one basis
\[
 \ket{x}=\frac{\ket{11}-\ket{00}}{\sqrt2},\quad
 \ket{y}=\frac{i(\ket{11}+\ket{00})}{\sqrt2},\quad
 \ket{z}=\frac{\ket{01}+\ket{10}}{\sqrt2},
\]
put $A=\id_3-T$ and let $[\mathbf r]_\times\mathbf v=
\mathbf r\times\mathbf v$. Substitution in \eqref{eq:pauli} gives
\begin{equation}\label{eq:cartesian}
 2\rho_2=A+i[\mathbf r]_\times,\qquad
 8\det\rho_2=\det A-\mathbf r^{\mathsf T}A\mathbf r.
\end{equation}
To check the determinant identity, diagonalize $A$ by a proper real
rotation. The matrix then has diagonal $a,b,c$ and off-diagonal entries
$-ir_3,ir_2,-ir_1$ and their conjugates; its determinant is
$abc-ar_1^2-br_2^2-cr_3^2$.

There is a useful bound specific to reductions of symmetric $N$-qubit
states. With $J_a=\frac12\sum_{n=1}^N\sigma_a^{(n)}$, permutation
symmetry gives
\[
 \langle J_a\rangle=\frac N2r_a,\qquad
 \left\langle\frac{J_aJ_b+J_bJ_a}{2}\right\rangle
 =\frac N4\delta_{ab}+\frac{N(N-1)}4T_{ab}.
\]
The nonnegative variance of every real linear combination of the $J_a$
implies $\id+(N-1)T-N\mathbf r\mathbf r^{\mathsf T}\succeq0$.
Testing \eqref{eq:cartesian} on real vectors gives $A\succeq0$.
Consequently,
\begin{equation}\label{eq:A-bounds}
 0\preceq A\preceq\frac{N}{N-1}\id_3\preceq\frac32\id_3,
 \qquad\Tr A=2.
\end{equation}
Only the uniform bound $A\preceq3\id_3/2$ is needed below.

Write $\rho'=\Phi(\rho)$. The partial trace intertwines rotations,
so the rank-one and rank-two parts of the reduction acquire the same
multipliers $f_1=f$ and $f_2=g$ as the global tensor components.
Equivalently, the above first and second moments give directly
\begin{equation}\label{eq:transform}
 \mathbf r'=f\mathbf r,\quad T'=gT+\frac{1-g}{3}\id_3,
 \quad A'=\alpha\id_3+gA,\qquad\alpha=\frac{2(1-g)}3.
\end{equation}

\section{Determinant monotonicity}
We prove two elementary estimates for $A$ and $A'$.

\begin{lemma}\label{lem:volume}
Under \eqref{eq:A-bounds} and \eqref{eq:region}, $\det A'\ge\det A$.
\end{lemma}
\begin{proof}
Write the eigenvalues of $A$ as $2/3+x_i$, where $\sum_i x_i=0$,
and set $s_2=\sum_i x_i^2$, $p=x_1x_2x_3$. Expansion gives
\begin{equation}\label{eq:volume-expansion}
 \det A'-\det A=\frac{1-g^2}{3}s_2-(1-g^3)p.
\end{equation}
We claim $p\le5s_2/36$. This is immediate if $p\le0$.
If $p>0$, exactly one $x_i$ is positive, so after relabeling
$(x_1,x_2,x_3)=(q+r,-q,-r)$ with $q,r>0$.
The upper eigenvalue bound gives $q+r\le5/6$, whereas
$s_2=2(q^2+qr+r^2)\ge6qr$. Hence
$p=(q+r)qr\le5s_2/36$, as claimed.
Since $g\le1$, \eqref{eq:volume-expansion} now gives
\[
 \det A'-\det A\ge
 \frac{(1-g)s_2}{36}(7+7g-5g^2)\ge0,
\]
because $7+7g-5g^2=1+(5g+3)(2-g)\ge1$ on $[-3/5,1]$.
\end{proof}

\begin{lemma}[Cofactor comparison]\label{lem:cofactor}
If $b,c$ are any two eigenvalues of $A$, and
$b'=\alpha+gb$, $c'=\alpha+gc$, then
\begin{equation}\label{eq:cofactor}
 b'c'\ge f^2bc.
\end{equation}
\end{lemma}
\begin{proof}
We have $0\le b,c\le3/2$ and $b+c\le2$. Also
$\alpha+gx\ge0$ for $0\le x\le3/2$: for $g\ge0$ this is
immediate, and for $g<0$ it follows from
\begin{equation}\label{eq:positive}
 \alpha+gx\ge\alpha+\frac32g=\frac{4+5g}{6}\ge\frac16.
\end{equation}
Write $s=b+c$, $p=bc$.

If $g\ge-1/2$, then $2g+\alpha=2(1+2g)/3\ge0$.
The identity
\[
 4\bigl[b'c'-(g+\alpha)^2bc\bigr]
 =\alpha\bigl[(2g+\alpha)(b-c)^2
       +(2-s)\{2\alpha+(2g+\alpha)s\}\bigr]\ge0
\]
and $(g+\alpha)^2=(2+g)^2/9\ge f^2$ prove the claim.

If $-3/5\le g<-1/2$, Lemma~\ref{lem:multipliers} gives $f^2<1/9$.
Set $h=g^2-1/9>0$. The inequality
$(3/2-b)(3/2-c)\ge0$ gives $p\ge3s/2-9/4$.
Furthermore
\[
 \alpha g+\frac32h=\frac{(5g-1)(g+1)}6<0.
\]
It follows that
\begin{align*}
 b'c'-\frac{bc}{9}
 &=\alpha^2+\alpha gs+hp\\
 &\ge\alpha^2-\frac94h+
       \left(\alpha g+\frac32h\right)s\\
 &\ge\alpha^2+2\alpha g+\frac34h
   =\frac{13+16g-5g^2}{36}\ge\frac{2}{45}>0.
\end{align*}
The last bound follows from
$13+16g-5g^2=8/5+(g+3/5)(19-5g)$.
This proves \eqref{eq:cofactor} in both cases.
\end{proof}

We can now show
\begin{equation}\label{eq:det-monotone}
 \det\rho'_2\ge\det\rho_2.
\end{equation}
In a common eigenbasis of $A,A'$, with respective eigenvalues $a,b,c$
and $a',b',c'$, all six eigenvalues are nonnegative. Set
$\Delta=8\det\rho_2$, $\Delta'=8\det\rho'_2$.
Equation~\eqref{eq:cartesian} gives the polynomial identity
\begin{align*}
 abc(\Delta'-\Delta)
 &=(a'b'c'-abc)\Delta\\
 &\quad+aa'r_1^2(b'c'-f^2bc)
       +bb'r_2^2(a'c'-f^2ac)
       +cc'r_3^2(a'b'-f^2ab)\ge0.
\end{align*}
Every term is nonnegative by Lemmas~\ref{lem:volume} and
\ref{lem:cofactor}, and $\Delta\ge0$ follows from $\rho_2\succeq0$.
If $abc>0$, division proves \eqref{eq:det-monotone}.
If $A$ is singular, then
\eqref{eq:cartesian}, $A\succeq0$ and $\det\rho_2\ge0$ force
$\det\rho_2=0$. The inequality is then immediate from
$\rho'_2\succeq0$. This also covers all boundary states.

\section{From determinant and purity to fidelity}
Orthogonally decompose the reduction into its tensor ranks:
\[
 \rho_2=\frac{\id_3}{3}+X_1+X_2,\qquad
 \rho'_2=\frac{\id_3}{3}+fX_1+gX_2.
\]
The Hermitian $X_1,X_2$ are traceless and Hilbert--Schmidt orthogonal.
The bounds $|f|,|g|\le1$ therefore imply
\begin{equation}\label{eq:purity}
 \Tr[(\rho'_2)^2]=\frac13+f^2\Tr(X_1^2)+g^2\Tr(X_2^2)
 \le\Tr(\rho_2^2).
\end{equation}
For a positive trace-one $3\times3$ matrix $B$, let
$e_2(B)=(1-\Tr B^2)/2$ and $e_3(B)=\det B$.
Equations~\eqref{eq:det-monotone} and \eqref{eq:purity} show that
both $e_2$ and $e_3$ increase. The following elementary fact completes
the argument without an integral representation of the square root.

\begin{lemma}[Three-eigenvalue comparison]\label{lem:roots}
If $B,C\succeq0$ are $3\times3$ matrices of trace one with
$e_2(C)\ge e_2(B)$ and $\det C\ge\det B$, then
$\Tr\sqrt C\ge\Tr\sqrt B$.
\end{lemma}
\begin{proof}
Let $x,y,z\ge0$ be the square roots of the eigenvalues of $B$.
Write $v=x+y+z$, $d=xyz=\sqrt{\det B}$ and $e=e_2(B)$.
Since $x^2+y^2+z^2=1$, elementary expansion gives
\begin{equation}\label{eq:root-identities}
 v\ge1,\qquad (v^2-1)^2=4e+8dv,\qquad v(v^2-1)\ge18d.
\end{equation}
For the last inequality, use $v^2-1=2(xy+xz+yz)$ and
\[
 v(xy+xz+yz)-9xyz
 =z(x-y)^2+x(y-z)^2+y(z-x)^2\ge0.
\]
Define $u,d',e'$ in the same way for $C$. If $u<v$, then
$u\ge1$, $d'\ge d$ and $e'\ge e$ imply
\[
 (v^2-1)^2-(u^2-1)^2
 =4(e-e')+8(dv-d'u)\le8d(v-u).
\]
Dividing by $v-u>0$ gives
$(v+u)(v^2+u^2-2)\le8d$. But
\[
 (v+u)(v^2+u^2-2)
 =v(v^2-1)+u(v^2-1)+(v+u)(u^2-1)
 >v(v^2-1)\ge18d\ge8d,
\]
a contradiction. Here strictness follows from $u\ge1$ and $v>u$.
\end{proof}

Apply Lemma~\ref{lem:roots} to $B=\rho_2$, $C=\rho'_2$.
Then $\Tr\sqrt{\rho'_2}\ge\Tr\sqrt{\rho_2}\ge0$; squaring and
using \eqref{eq:candidate} proves \eqref{eq:main}.
This proves precisely the fidelity part of Conjecture~1 at $t=2$.
\qed

\section{Formal verification and reproducibility}
The accompanying Lean~4.24.0 files certify the complete scalar algebraic
reduction of the argument, including every $N\ge3$, arbitrary probability
mixtures, and singular inputs. The main theorem is
\texttt{SU2T2.channel\_fidelity\_monotone\_from\_scalar\_reduction}.
Its explicit hypotheses encode the multiplier formulas, the marginal
constraints, and the spectral purity and determinant identities.
The files contain no unfinished proofs or additional axioms. Their exact
scope, dependency pins, checked-source hashes, and build instructions
are recorded in \texttt{lean/README.md}.

This is not an end-to-end Lean formalization of Theorem~\ref{thm:main}:
the Choi/Casimir reduction, the physical marginal constraints, and the
identification with spectral data are proved above but are not themselves
formalized in Lean. This distinction is part of the certificate's scope.

\begin{remark}[Why the marginal constraint matters]
The bound $A\preceq3\id_3/2$ comes from the symmetric extension to
$N\ge3$ qubits. An arbitrary spin-one density matrix need not obey it.
The proof therefore does not assume that the induced reduced map is a
positive or completely positive map on the whole spin-one state space.
\end{remark}
\end{document}
